\documentclass[11pt,a4paper]{article}
\usepackage[utf8]{inputenc}
\usepackage[T1]{fontenc}
\usepackage[spanish]{babel}
\usepackage{amsmath,amssymb,mathtools,bm}
\usepackage{booktabs}
\usepackage{geometry}
\usepackage{hyperref}
\usepackage{siunitx}
\usepackage{enumitem}
\geometry{margin=2.2cm}

\title{Guía Didáctica de la Implementación Actual\\
PROM--RBF, ECM y MAW--ECM (2D)}
\author{Basada en la implementación actual del proyecto\\
\texttt{RVE\_homogenization\_NeoHookean\_using\_Kratos\_2D\_MAWECM}}
\date{\today}

\begin{document}
\maketitle
\tableofcontents

\section{Objetivo de esta guía}
Esta guía está escrita para alguien que \textbf{empieza desde cero} y quiere entender:
\begin{itemize}[leftmargin=1.5em]
\item qué significa cada símbolo;
\item qué tamaño tiene cada matriz/vector;
\item cómo se conectan los \textit{stages} 0--8;
\item qué hace exactamente ECM clásico;
\item qué hace exactamente MAW--ECM;
\item qué cambia en el modo nuevo \texttt{single\_res\_hom}.
\end{itemize}

\noindent
La idea es que puedas explicarlo a otro estudiante de máster paso a paso, sin saltos.

\section{Mapa mental completo (de Stage0 a Stage8)}
\subsection{Pipeline conceptual}
\begin{enumerate}[leftmargin=1.5em]
\item \textbf{Stage0}: defines malla estructurada en espacio paramétrico $\mu$ y dos trayectorias.
\item \textbf{Stage1}: corres FOM en esas trayectorias y guardas snapshots.
\item \textbf{Stage2a}: construyes POD en DOFs libres (\textit{fluctuaciones}).
\item \textbf{Stage2b}: construyes coordenadas maestro-esclavo (LS + descomposición ortogonal).
\item \textbf{Stage3}: armas dataset final de regresión y topología estructurada.
\item \textbf{Stage4}: entrenas RBF para el decoder $q_m \mapsto q_s$ (o variantes).
\item \textbf{Stage5}: validas PROM-RBF online contra FOM.
\item \textbf{Stage6}: construyes ECM clásico (residual y homogenización).
\item \textbf{Stage7}: validas HPROM-RBF clásico.
\item \textbf{Stage8a}: construyes dataset estructurado MAW (residual, y homogenización opcional).
\item \textbf{Stage8b}: entrenas MAW-ECM (fase actual: pruning local sin grafo; modo residual o single residual+hom).
\item \textbf{Stage8 test}: validas HPROM-MAWECM-RBF online.
\end{enumerate}

\subsection{Números reales de tu corrida actual (2D)}
En una corrida típica tuya:
\begin{itemize}[leftmargin=1.5em]
\item DOFs totales: $4244$.
\item DOFs libres: $3924$.
\item DOFs Dirichlet: $320$.
\item snapshots FOM: $5574$.
\item dimensión POD: $r=13$.
\item dimensión master: $d_m=2$.
\item dimensión slave: $d_s=11$.
\item elementos FE: $n_e=990$.
\item nodos estructurados Stage0/MAW: $N_n=451$.
\end{itemize}

\section{Diccionario de símbolos (con tamaño)}
\subsection{Convención de notación (para no colisionar símbolos)}
En esta guía voy a usar \textbf{siempre} esta convención:
\begin{itemize}[leftmargin=1.5em]
\item \textbf{Pesos ECM/MAW}: $w$ (y matrices $W_{train}$ cuando hay muchos nodos).
\item \textbf{Matriz de snapshots}: $X_{snap}$ (nunca $W$).
\item \textbf{Adyacencia de grafo}: $A_{adj}$ (nunca $W$).
\item \textbf{Diagonal de grados}: $D_{deg}$ (nunca $D$ a secas).
\item \textbf{Área de elemento}: $|\Omega_e|$ (no $A_e$ para evitar choque con otras $A$).
\item \textbf{Matrices tipo ``A''}: siempre con subíndice explícito:
  $A_m$ (master), $A_j$ (restricción local MAW), $A_{adj}$ (grafo).
\end{itemize}

\subsection{Campos físicos y discretización}
\begin{itemize}[leftmargin=1.5em]
\item $\mathbf{u}\in\mathbb{R}^{n_{dof}}$: desplazamientos globales.
\item $\mathbf{u}_f\in\mathbb{R}^{n_f}$: desplazamientos en DOFs libres.
\item $\mathbf{u}_d\in\mathbb{R}^{n_d}$: desplazamientos en DOFs fijos (Dirichlet).
\item $n_{dof}=n_f+n_d$.
\item $\mathbf{R}_f(\mathbf{u})\in\mathbb{R}^{n_f}$: residual en espacio libre.
\item $\mathbf{K}_{ff}(\mathbf{u})\in\mathbb{R}^{n_f\times n_f}$: tangente libre-libre.
\end{itemize}

\subsection{POD y coordenadas reducidas}
\begin{itemize}[leftmargin=1.5em]
\item $\Phi_{rom}\in\mathbb{R}^{n_f\times r}$: base POD total.
\item $\mathbf{q}_{pod}\in\mathbb{R}^{r}$: coordenadas POD.
\item $\Phi_m\in\mathbb{R}^{n_f\times d_m}$: base master en desplazamientos.
\item $\Phi_s\in\mathbb{R}^{n_f\times d_s}$: base slave en desplazamientos.
\item $\mathbf{q}_m\in\mathbb{R}^{d_m}$: coordenadas master.
\item $\mathbf{q}_s\in\mathbb{R}^{d_s}$: coordenadas slave.
\item $d_m+d_s=r$ (en tu caso $2+11=13$).
\end{itemize}

\subsection{Operadores LS/master-slave}
\begin{itemize}[leftmargin=1.5em]
\item $\mathbf{T}_m\in\mathbb{R}^{d_m\times r}$: mapa LS desde $q_{pod}$ a objetivo master.
\item $\mathbf{A}_m\in\mathbb{R}^{d_m\times d_m}$: factor interno de la base master, obtenido de la SVD de $\Phi_c=\Phi_{rom}T_m^\top$ (sección Stage2b).
\item $\mathbf{C}_m\in\mathbb{R}^{r\times d_m}$: base master en espacio POD.
\item $\mathbf{C}_s\in\mathbb{R}^{r\times d_s}$: base slave en espacio POD.
\end{itemize}

\subsection{ECM/MAW}
\begin{itemize}[leftmargin=1.5em]
\item $n_e$: número de elementos.
\item $\mathcal{Z}$: conjunto de elementos seleccionados.
\item $\mathbf{w}\in\mathbb{R}^{|\mathcal{Z}|}$: pesos ECM en soporte seleccionado.
\item $\mathbf{w}^{full}\in\mathbb{R}^{n_e}$: mismos pesos embebidos en malla completa.
\item $\mathbf{Q}_{ecm}\in\mathbb{R}^{(n_q N_s^{res})\times n_e}$: bloques residual proyectado.
\item $\mathbf{b}_{res}\in\mathbb{R}^{n_q N_s^{res}}$: suma completa residual.
\item $\mathbf{C}_{hom}\in\mathbb{R}^{(6N_s^{hom})\times n_e}$: bloques homogenización.
\item $\mathbf{b}_{hom}\in\mathbb{R}^{6N_s^{hom}}$: suma completa homogenización.
\item $A_{adj}\in\mathbb{R}^{N_n\times N_n}$: adyacencia del grafo estructurado.
\item $D_{deg}\in\mathbb{R}^{N_n\times N_n}$: diagonal de grados del grafo estructurado.
\item $K_{graph}=D_{deg}-A_{adj}\in\mathbb{R}^{N_n\times N_n}$: Laplaciano usado en MAW.
\item $A_j\in\mathbb{R}^{m_j\times n_c}$: operador local MAW en nodo estructurado $j$.
\item $b_j\in\mathbb{R}^{m_j}$: lado derecho local MAW en nodo estructurado $j$.
\end{itemize}

\section{Stage0: espacio paramétrico y trayectorias estructuradas}
\subsection{Qué es el espacio paramétrico aquí}
En esta implementación, el espacio paramétrico base es:
\[
\mu = [G_x,\; G_{xy}]
\]
donde $G_x$ y $G_{xy}$ parametrizan una deformación macroscópica 2D (en este caso, rama triangular superior de $F$ por construcción del script).

\subsection{Malla estructurada}
Se crea una malla cartesiana en $\mu$:
\[
G_x \in [G_x^{min}, G_x^{max}],\qquad
G_{xy} \in [G_{xy}^{min}, G_{xy}^{max}]
\]
con $n_{gx}\times n_{gxy}$ nodos.

En tu configuración usual:
\[
n_{gx}=41,\quad n_{gxy}=11,\quad N_n=41\cdot 11=451.
\]

\subsection{Grafo y Laplaciano}
Se construye un grafo de vecindad sobre esa malla:
\begin{itemize}[leftmargin=1.5em]
\item nodo = punto estructurado en $\mu$;
\item arista = vecino horizontal o vertical.
\end{itemize}
Con eso:
\[
K_{graph} = D_{deg} - A_{adj},
\]
donde $A_{adj}$ es adyacencia y $D_{deg}$ es diagonal de grados.

\textbf{Importante:} este grafo se guarda para análisis y para una fase 2 con acoplamiento espacial.
En la fase 1 actual de MAW, el pruning se ejecuta sin acoplamiento por grafo.

\subsection{Mapeo de $\mu$ a deformación objetivo de Stage1}
Se guardan waypoints de strain $[E_{xx},E_{yy},\gamma_{xy}]$ usando:
\begin{align}
E_{xx} &= G_x + \tfrac12 G_x^2,\\
E_{yy} &= \tfrac12 G_{xy}^2,\\
\gamma_{xy} &= G_{xy}(1+G_x),
\end{align}
cuando \texttt{mapping=green\_lagrange\_upper}.

\section{Stage1: FOM}
Se recorren trayectorias y se resuelve el problema no lineal completo.
Se guarda por snapshot:
\begin{itemize}[leftmargin=1.5em]
\item $\mathbf{u}$ completa;
\item strain aplicado;
\item variables de homogenización.
\end{itemize}

\textbf{Detalle clave}: el número de incrementos por trayectoria está escalado por \texttt{ref\_steps} (por ejemplo 400), lo que controla robustez de convergencia.

\section{Stage2a: POD de fluctuaciones}
\subsection{Separación affine + fluctuación}
Para cada snapshot:
\[
\mathbf{u}_f = \mathbf{u}_{aff,f}(E) + \mathbf{u}_{fluc},
\]
donde $\mathbf{u}_{aff,f}(E)$ se calcula desde la parte macroscópica y
$\mathbf{u}_{fluc}$ es la fluctuación.

\subsection{Construcción POD}
Se apilan fluctuaciones:
\[
X_{snap} \in \mathbb{R}^{N_s \times n_f},\quad X_{snap}^\top\in\mathbb{R}^{n_f\times N_s}.
\]
Se hace SVD:
\[
X_{snap}^\top = U\Sigma V^\top,
\]
y se toma $\Phi_{rom}=U_{(:,1:r)}$.

Entonces:
\[
\mathbf{q}_{pod} = \Phi_{rom}^\top \mathbf{u}_{fluc},
\quad\text{(equivalente por filas a }Q_{pod}=X_{snap}\Phi_{rom}\text{)}.
\]

\textbf{En tu caso}: $r=13$.

\section{Stage2b: LS master + descomposición maestro/esclavo}
\subsection{Paso 1: recuperar $\mu$ desde strain aplicado}
Desde $E_{xx},\gamma_{xy}$ se invierte el mapeo para obtener
$[G_x,G_{xy}]$ (o $[F_{11},F_{12}]$).

\subsection{Paso 2: ajuste LS $q_{pod}\to \mu$}
Con muestras:
\[
Q\in\mathbb{R}^{N_s\times r},\qquad
M\in\mathbb{R}^{N_s\times d_m}\;(d_m=2).
\]
Se resuelve:
\[
\min_{T_m}\|Q\,T_m^\top - M\|_F^2,
\]
guardando $T_m$.

\subsubsection*{¿Cómo se resuelve \texorpdfstring{$\min_{T_m}$}{min Tm} en el código?}
En \texttt{stage2b\_build\_ls\_master\_from\_pod.py} se usa:
\[
\texttt{np.linalg.lstsq}(Q, M,\texttt{rcond=None}),
\]
que computa una solución de mínimos cuadrados por SVD/pseudo-inversa.
En forma compacta:
\[
T_m^\top = Q^+ M,
\]
donde $Q^+$ es la pseudo-inversa de Moore--Penrose.
Si $Q$ tuviera rango columna completo, esto coincide con:
\[
T_m^\top = (Q^\top Q)^{-1}Q^\top M.
\]
\textbf{Importante:} aquí no se hace una regularización explícita tipo ridge; la robustez viene del solver SVD interno de \texttt{lstsq}.

\subsection{Paso 3: base master en espacio físico}
Con $T_m^\top\in\mathbb{R}^{r\times d_m}$:
\[
\Phi_c=\Phi_{rom}T_m^\top\in\mathbb{R}^{n_f\times d_m}.
\]
SVD de $\Phi_c$:
\[
\Phi_c = U_m\Sigma_m V_m^\top.
\]
\[
\Phi_m:=U_m(:,1{:}d_m),\qquad
A_m:=\Sigma_m V_m^\top.
\]
Así, por construcción:
\[
\Phi_m\in\mathbb{R}^{n_f\times d_m},\qquad
A_m\in\mathbb{R}^{d_m\times d_m}.
\]
\textbf{De aquí sale exactamente $A_m$}. No es un parámetro libre: es el factor derecho de la SVD truncada de $\Phi_c$.

\subsection{Paso 4: subespacio slave en POD}
Se define:
\[
C_m=\Phi_{rom}^\top\Phi_m\in\mathbb{R}^{r\times d_m}.
\]
Se construye su complemento ortonormal:
\[
C_s\in\mathbb{R}^{r\times d_s},\qquad d_s=r-d_m.
\]
\subsubsection*{¿De dónde sale exactamente \texorpdfstring{$C_s$}{Cs}? (matemática + código)}
La idea es: $C_s$ debe generar las direcciones POD que \textbf{no} están en el subespacio master.

\textbf{Paso A:} tomas $C_m$ y formas el proyector ortogonal al subespacio master:
\[
P_\perp = I_r - C_m C_m^\top \in \mathbb{R}^{r\times r}.
\]
Si $C_m$ es ortonormal por columnas, entonces:
\[
C_m^\top C_m = I_{d_m},
\]
y $P_\perp$ proyecta a la parte ``slave'' del espacio POD.

\textbf{Paso B:} calculas autodescomposición de $P_\perp$:
\[
P_\perp v_i = \lambda_i v_i.
\]
Los autovectores con autovalores cercanos a $1$ span-ean el complemento ortogonal.
Como $d_s=r-d_m$, tomas exactamente $d_s$ vectores:
\[
C_s = [v_1,\dots,v_{d_s}] \in \mathbb{R}^{r\times d_s}.
\]

\textbf{Paso C (numérico):} re-ortonormalizas $C_s$ con QR para robustez:
\[
C_s \leftarrow \operatorname{qr}(C_s),
\]
de modo que en la práctica se cumpla:
\[
C_s^\top C_s \approx I_{d_s},\qquad C_m^\top C_s \approx 0.
\]

\textbf{Esto coincide con el código} de \texttt{stage2b\_build\_ls\_master\_from\_pod.py}:
\begin{itemize}[leftmargin=1.5em]
\item construye \texttt{proj\_perp = I - C\_m @ C\_m.T};
\item usa \texttt{np.linalg.eigh(proj\_perp)};
\item ordena autovectores y toma los primeros $d_s$;
\item aplica \texttt{np.linalg.qr} para ortonormalizar.
\end{itemize}

\textbf{Interpretación práctica:}
\begin{itemize}[leftmargin=1.5em]
\item $q_{pod}C_m$ = coordenadas de la parte master dentro del POD;
\item $q_{pod}C_s$ = coordenadas de la parte slave dentro del POD;
\item juntas reconstruyen las fluctuaciones POD sin perder información del subespacio reducido de rango $r$.
\end{itemize}

Para cada snapshot:
\[
q_{m,proj}=q_{pod}C_m,\qquad q_s=q_{pod}C_s.
\]

Y para consistencia exacta con $A_m$:
\[
q_m^{cons} = A_m^{-1}q_{m,proj}.
\]

\textbf{Muy importante (tu duda recurrente):}
\begin{itemize}[leftmargin=1.5em]
\item el $q_m$ usado en Stage3+online es $q_m^{cons}$ (no el LS ``crudo'');
\item por eso el decoder queda coherente con
\(
q_{master}=A_m q_m
\)
y con la reconstrucción.
\end{itemize}

\subsection{Mini ejemplo numérico (didáctico)}
Supón:
\[
r=4,\quad d_m=2,\quad d_s=2.
\]
Para un snapshot:
\[
q_{pod}=[0.8,\,-0.2,\,0.3,\,1.1].
\]
Si
\[
C_m=
\begin{bmatrix}
1&0\\
0&1\\
0&0\\
0&0
\end{bmatrix},
\quad
C_s=
\begin{bmatrix}
0&0\\
0&0\\
1&0\\
0&1
\end{bmatrix},
\]
entonces:
\[
q_{m,proj}=q_{pod}C_m=[0.8,\,-0.2],\quad
q_s=q_{pod}C_s=[0.3,\,1.1].
\]
Si además
\[
A_m=
\begin{bmatrix}
2&0\\
0&1
\end{bmatrix},
\]
entonces:
\[
q_m^{cons}=A_m^{-1}q_{m,proj}=[0.4,\,-0.2].
\]
Y se cumple:
\[
A_m q_m^{cons}=q_{m,proj}.
\]

\section{Stage3: dataset de regresión}
Se empaqueta:
\begin{itemize}[leftmargin=1.5em]
\item entradas $\mu$;
\item targets (en tu configuración actual: $q_s$);
\item arrays completos $q_m,q_s,q_{pod}$;
\item matrices $A_m,C_m,C_s$;
\item topología estructurada Stage0 (nodos, quads, grafo, Laplaciano);
\item mapeo nodo estructurado $\leftrightarrow$ snapshot más cercano.
\end{itemize}

\textbf{Tu corrida}: $5574$ muestras; split $80/10/10$.

\section{Stage4: RBF para decoder}
\subsection{Qué aprende}
En el flujo actual, se usa:
\[
\text{input} = q_m \in \mathbb{R}^{2},
\qquad
\text{target} = q_s \in \mathbb{R}^{11}.
\]
O sea, aprende:
\[
N_{slave}: q_m \mapsto q_s.
\]

\subsection{Modelo RBF}
Con centros $c_\ell$:
\[
\phi_\ell(q)=\exp\!\left(-\tfrac12\sum_{a=1}^{d_m}\frac{(q_a-c_{\ell a})^2}{\ell_a^2}\right).
\]
La salida es multi-output (11 componentes), con regularización y posible término polinomial.

\subsection{Por qué ``sparse''}
No se usan necesariamente todos los puntos de entrenamiento como centros.
Se seleccionan puntos inducidos (\textit{inducing points}) para balancear costo/precisión.

\section{PROM online: ecuaciones que realmente se resuelven}
En Stage5/7/8, el solver PROM-RBF no ``solo predice'': hace corrector de Newton en $q_m$.

\subsection{Ansatz de desplazamiento}
\[
\mathbf{u}_f(q_m)=\mathbf{u}_{aff,f}(E)+\Phi_m(A_m q_m)+\Phi_s q_s(q_m).
\]

\subsection{Jacobiano cinemático reducido}
\[
\frac{\partial \mathbf{u}_f}{\partial q_m}
=
\Phi_m A_m
+
\Phi_s\frac{\partial q_s}{\partial q_m},
\]
donde $\partial q_s/\partial q_m$ se calcula de forma \textbf{analítica} a partir del modelo RBF
(derivando kernel y parte polinómica local; en modo \texttt{neighbors}, usando el vecindario activo).

\subsection{Residual y tangente reducidos}
\[
r_r(q_m)=
\left(\frac{\partial \mathbf{u}_f}{\partial q_m}\right)^\top
R_f(\mathbf{u}_f(q_m)),
\]
\[
K_r(q_m)=
\left(\frac{\partial \mathbf{u}_f}{\partial q_m}\right)^\top
K_{ff}(\mathbf{u}_f(q_m))
\left(\frac{\partial \mathbf{u}_f}{\partial q_m}\right).
\]

\subsection{Paso de Newton reducido}
\[
K_r \Delta q_m = r_r,
\qquad
q_m \leftarrow q_m + \Delta q_m.
\]

\textbf{Nota de convergencia clave}: se monitorea $\|R_r\|$ (no $\|R_f\|$) para criterio reducido.
Esto explica por qué puede verse $R_f$ grande pero solver reducido estable.

\section{Stage6a: construcción de matrices ECM clásicas}
\subsection{Qué se construye exactamente}
En Stage6a no se ``entrena'' aún el ECM. Aquí solo se construyen las matrices
que definen el problema de cubatura.

\begin{itemize}[leftmargin=1.5em]
\item Entrada: contribuciones por elemento en muchos estados de entrenamiento.
\item Salida: matrices densas por bloques:
  \[
  Q_{res},\; b_{res},\; C_{hom},\; b_{hom}.
  \]
\item Estas matrices quedan guardadas en:
\texttt{Q\_ecm.dat}, \texttt{b\_full.dat}, \texttt{C\_hom.dat}, \texttt{b\_hom.dat}.
\end{itemize}

\subsection{Bloque residual: de dónde sale $Q_{res}$}
Para cada estado (snapshot) $s=1,\dots,N_s^{res}$ y elemento
$e=1,\dots,n_e$, se calcula un vector
\[
q_{s,e}\in\mathbb{R}^{n_q},
\]
donde $n_q$ es el número de ecuaciones reducidas que quieres reproducir
(en tu caso típico, $n_q=13$).

La interpretación de $q_{s,e}$ es:
\begin{itemize}[leftmargin=1.5em]
\item componente a componente, cuánto aporta \textbf{el elemento $e$}
a cada ecuación reducida del residual en el estado $s$.
\end{itemize}

\subsection{Receta exacta: del residual elemental $r_e$ al vector proyectado $q_{s,e}$}
Esta es la parte que normalmente queda implícita, pero aquí va explícita tal como
está implementada en \texttt{stage6a\_build\_ecm\_dataset.py}.

\subsubsection*{Paso A: residual elemental local}
Para un estado $s$ y elemento $e$, Kratos devuelve:
\[
r_e^{(s)}\in\mathbb{R}^{n_{loc}(e)},
\]
con
\[
r_e^{(s)}=\texttt{elem.CalculateRightHandSide(...)}.
\]
Este vector está en el orden local de DOFs del elemento.

\subsubsection*{Paso B: quedarse solo con DOFs que realmente participan}
El elemento también da su vector de IDs de ecuación local:
\[
\mathrm{id}_e \in \mathbb{Z}^{n_{loc}(e)},
\quad
\mathrm{id}_{e,a}=\texttt{EquationIdVector}[a].
\]
\begin{itemize}[leftmargin=1.5em]
\item Si $\mathrm{id}_{e,a}<0$, ese DOF no participa en el sistema lineal global $\Rightarrow$ se descarta.
\item Si $\mathrm{id}_{e,a}\ge 0$ pero es DOF de Dirichlet (no libre), también se descarta para la proyección residual.
\end{itemize}

Define entonces el conjunto de índices locales válidos:
\[
\mathcal I_{e}^{free}
=
\{a\in\{1,\dots,n_{loc}(e)\}\;:\;\mathrm{id}_{e,a}\ge 0,\;\text{dof libre}\}.
\]
Con eso:
\[
r_{e,free}^{(s)}\in\mathbb{R}^{n_{ef}},\qquad
n_{ef}=|\mathcal I_{e}^{free}|,
\]
tomando las entradas de $r_e^{(s)}$ en $\mathcal I_{e}^{free}$.

\subsubsection*{Paso C: proyector reducido local del elemento}
Sea
\[
\Phi_f\in\mathbb{R}^{n_f\times n_q}
\]
la base POD de DOFs libres (archivo \texttt{pod\_basis\_free.npy}).

Para cada índice local $a\in\mathcal I_{e}^{free}$, su DOF global libre asociado
es $p(a)\in\{1,\dots,n_f\}$ (en código: \texttt{map\_g2f}).

Construimos la submatriz:
\[
\Phi_{e,free}
\in\mathbb{R}^{n_{ef}\times n_q},
\qquad
\Phi_{e,free}(a',i)=\Phi_f\big(p(a'),i\big),
\]
y su traspuesta:
\[
\Phi_{e,free}^\top \in\mathbb{R}^{n_q\times n_{ef}}.
\]

\subsubsection*{Paso D: proyección residual por elemento}
La contribución proyectada del elemento $e$ en el estado $s$ es:
\[
q_{s,e}=
\Phi_{e,free}^\top r_{e,free}^{(s)}
\in\mathbb{R}^{n_q}.
\]
Componente a componente:
\[
q_{s,e}^{(i)}
=
\sum_{a'=1}^{n_{ef}}
\Phi_{e,free}(a',i)\,r_{e,free,a'}^{(s)}.
\]

\textbf{Esto es exactamente el ``residuo proyectado por elemento''} que luego entra a ECM.

\subsubsection*{Forma equivalente con operadores de restricción}
Si defines una matriz de selección local-global
\[
R_e\in\{0,1\}^{n_{ef}\times n_f},
\]
entonces
\[
\Phi_{e,free}=R_e\Phi_f,\qquad
q_{s,e}=(R_e\Phi_f)^\top r_{e,free}^{(s)}.
\]

\subsubsection*{Cómo se arma la fila-bloque de $Q_{res}$}
Para un estado $s$:
\[
Q_{res}^{(s)}=
\begin{bmatrix}
q_{s,1} & q_{s,2} & \cdots & q_{s,n_e}
\end{bmatrix}
\in\mathbb{R}^{n_q\times n_e}.
\]
Luego se apila verticalmente para todos los $s$:
\[
Q_{res}=
\begin{bmatrix}
Q_{res}^{(1)}\\
Q_{res}^{(2)}\\
\vdots\\
Q_{res}^{(N_s^{res})}
\end{bmatrix}
\in\mathbb{R}^{(n_qN_s^{res})\times n_e}.
\]

\subsubsection*{Qué representa $b_{res}$}
En cada estado:
\[
b_{res}^{(s)}=\sum_{e=1}^{n_e} q_{s,e}=Q_{res}^{(s)}\mathbf{1}.
\]
Globalmente:
\[
b_{res}=Q_{res}\mathbf{1}.
\]
Este $b_{res}$ es la referencia ``full integration'' que ECM intenta reproducir con pocos elementos y pesos.

\subsubsection*{Nota importante sobre esta implementación}
En este pipeline, la proyección de Stage6a usa la base fija $\Phi_f$.
Es decir, el operador de proyección offline es fijo en todos los snapshots.
No se está usando aquí un Jacobiano de variedad dependiente del estado para construir $Q_{res}$.
Eso es coherente con el código actual de Stage6a/Stage8a.

\subsection{Mini ejemplo de llenado (residual)}
Supón $N_s^{res}=2$, $n_q=2$, $n_e=3$.
Entonces $Q_{res}$ tiene tamaño $4\times3$.
Si
\[
q_{1,1}=\begin{bmatrix}1\\2\end{bmatrix},\;
q_{1,2}=\begin{bmatrix}3\\4\end{bmatrix},\;
q_{1,3}=\begin{bmatrix}5\\6\end{bmatrix},
\]
\[
q_{2,1}=\begin{bmatrix}7\\8\end{bmatrix},\;
q_{2,2}=\begin{bmatrix}9\\10\end{bmatrix},\;
q_{2,3}=\begin{bmatrix}11\\12\end{bmatrix},
\]
entonces
\[
Q_{res}=
\begin{bmatrix}
1&3&5\\
2&4&6\\
7&9&11\\
8&10&12
\end{bmatrix},
\qquad
b_{res}=Q_{res}\mathbf{1}=
\begin{bmatrix}
9\\12\\27\\30
\end{bmatrix}.
\]

\subsection{Bloque homogenización: de dónde sale $C_{hom}$}
Para cada estado $s$ y elemento $e$ se forma un vector de 6 componentes:
\[
c_{s,e}=
\begin{bmatrix}
|\Omega_e|\bar\varepsilon_{xx}\\
|\Omega_e|\bar\varepsilon_{yy}\\
|\Omega_e|\bar\gamma_{xy}\\
|\Omega_e|\bar\sigma_{xx}\\
|\Omega_e|\bar\sigma_{yy}\\
|\Omega_e|\bar\sigma_{xy}
\end{bmatrix}\in\mathbb{R}^{6}.
\]

\textbf{Muy importante:} cada componente ya está ponderada por medida del elemento
($|\Omega_e|$ en 2D). Por eso, al sumar por elementos, obtienes la contribución macroscópica.

Se arma:
\[
C_{hom}\in\mathbb{R}^{(6N_s^{hom})\times n_e},\qquad
b_{hom}=C_{hom}\mathbf{1}.
\]
Luego este bloque se separa en
\[
C_{\varepsilon}\in\mathbb{R}^{(3N_s^{hom})\times n_e},\qquad
C_{\sigma}\in\mathbb{R}^{(3N_s^{hom})\times n_e}
\]
para tratar strain y stress por separado cuando conviene.

\section{Stage6b: ECM clásico (independent/cascade/single)}
\subsection{Problema matemático ideal (lo que queremos)}
Dados $A\in\mathbb{R}^{m\times n_e}$ y $b\in\mathbb{R}^m$, ECM busca:
\[
\min_{w\in\mathbb{R}^{n_e}} \|w\|_0
\quad\text{sujeto a}\quad
\|Aw-b\|_2\le \tau\|b\|_2,\;\; w\ge 0.
\]
\begin{itemize}[leftmargin=1.5em]
\item $w$: pesos de cubatura por elemento.
\item $\|w\|_0$: número de elementos activos (sparseza).
\item $w\ge 0$: no negatividad (propiedad física/numérica deseada).
\end{itemize}

\subsection{Problema práctico (cómo lo resolvemos en código)}
Como el problema ideal es combinatorio, Stage6b usa esta secuencia:
\begin{enumerate}[leftmargin=1.5em]
\item \textbf{Compresión algebraica} de $A^\top$ con RSVD:
\[
A^\top \approx U\Sigma V^\top,
\]
y se retiene una base $U$ truncada por tolerancia.
\item \textbf{Greedy ECM} sobre esa base para seleccionar soporte
$\mathcal Z$ y pesos positivos.
\item \textbf{Reconstrucción full}: se crea $w_{full}\in\mathbb{R}^{n_e}$,
poniendo ceros fuera de $\mathcal Z$.
\end{enumerate}

\subsection{Cómo se resuelve numéricamente (funciones concretas)}
\subsubsection*{RSVD (en código)}
En \texttt{stage6b\_compute\_ecm\_weights.py} se llama:
\[
\texttt{RandomizedSingularValueDecomposition.Calculate}(A^\top,\texttt{truncation\_tolerance}=\tau_{svd}).
\]
Se conserva la base $U$ truncada y esa base es la entrada de ECM.

\subsubsection*{ECM clásico (en código)}
Se instancia:
\[
\texttt{EmpiricalCubatureMethod(ECM\_tolerance,\dots)}
\]
y luego:
\[
\texttt{SetUp(ResidualsBasis=U,\;constrain\_sum\_of\_weights=True,\dots)}.
\]
Ese \texttt{constrain\_sum\_of\_weights=True} añade una restricción auxiliar
de ``conteo'' para evitar la solución trivial y fijar una normalización global de pesos.

\subsubsection*{Paso greedy interno (detalle algebraico)}
En el ECM clásico (\texttt{empirical\_cubature\_method.py}):
\begin{itemize}[leftmargin=1.5em]
\item selección de candidato:
\[
i=\arg\max_{k\in\mathcal Y} g_k^\top r,
\]
donde $g_k$ es la columna candidata y $r$ es residual del problema ECM;
\item en la primera iteración:
\[
\alpha=\texttt{np.linalg.lstsq}(G_{(:,i)},b),
\]
\item en iteraciones siguientes, actualización rápida del inverso (Sherman-Morrison-like)
con \texttt{\_UpdateWeightsInverse}, sin recomputar todo desde cero;
\item si aparecen pesos negativos, se corrige active-set:
se sacan esos índices del soporte y se reevalúan pesos.
\end{itemize}

\subsubsection*{Reconstrucción final}
Si $\mathcal Z$ es soporte seleccionado y $w_{\mathcal Z}$ sus pesos:
\[
w_{full}(e)=
\begin{cases}
w_{\mathcal Z}(k), & e=\mathcal Z_k,\\
0, & e\notin\mathcal Z.
\end{cases}
\]

\subsection{Qué son exactamente $A$ y $b$ en cada modo}
\subsubsection*{Modo independent}
Tres problemas separados:
\[
A_{res}=Q_{res},\; b_{res};
\qquad
A_{\varepsilon}=C_{\varepsilon},\; b_{\varepsilon};
\qquad
A_{\sigma}=C_{\sigma},\; b_{\sigma}.
\]
Resultado: $(\mathcal Z_{res},w_{res})$, $(\mathcal Z_{\varepsilon},w_{\varepsilon})$,
$(\mathcal Z_{\sigma},w_{\sigma})$.

\subsubsection*{Modo cascade}
Mismos $A,b$ que \texttt{independent}, pero inicializando cada selección con el soporte previo
($\mathcal Z_{res}\to\mathcal Z_{\varepsilon}\to\mathcal Z_{\sigma}$) para favorecer consistencia.

\subsubsection*{Modo single}
Un solo problema combinado:
\[
\widetilde{A}=
\begin{bmatrix}
\alpha_{res}Q_{res}\\
\alpha_{\varepsilon}C_{\varepsilon}\\
\alpha_{\sigma}C_{\sigma}
\end{bmatrix},
\qquad
\widetilde{b}=
\begin{bmatrix}
\alpha_{res}b_{res}\\
\alpha_{\varepsilon}b_{\varepsilon}\\
\alpha_{\sigma}b_{\sigma}
\end{bmatrix}.
\]
Los factores $\alpha$ normalizan bloques (por norma Frobenius, por fila o sin normalizar)
para que un bloque no domine numéricamente al resto.

\subsection{Qué sale en \texttt{ecm\_weights\_all.npz}}
Las claves más importantes para entender el pipeline son:
\[
Z_{res},Z_{\varepsilon},Z_{\sigma},Z_{union},
\quad
w_{res},w_{\varepsilon},w_{\sigma},
\quad
w_{res}^{full},w_{\varepsilon}^{full},w_{\sigma}^{full}.
\]
\begin{itemize}[leftmargin=1.5em]
\item $Z_\star$: índices de elementos seleccionados.
\item $w_\star$: pesos solo del soporte.
\item $w_\star^{full}$: vector largo de tamaño $n_e$ con ceros fuera del soporte.
\end{itemize}

\section{Stage6c: creación de malla HROM (\texttt{.mdpa})}
Aquí pasas de ``soporte algebraico'' a ``malla física''.

\subsection{Qué hace}
\begin{enumerate}[leftmargin=1.5em]
\item Toma un conjunto de selección (por ejemplo \texttt{Z\_union}).
\item Extrae nodos/elementos necesarios del \texttt{.mdpa} completo.
\item Escribe un \texttt{.mdpa} reducido para corridas HROM.
\item Guarda metadatos de mapeo full$\leftrightarrow$hrom.
\end{enumerate}

\subsection{Por qué es clave}
Sin este paso, podrías tener pesos ECM correctos pero aún estar montando
sistemas sobre la malla completa en el solver.
Con este paso, el ensamblaje también se restringe a los elementos seleccionados.

\section{Stage7: HPROM-RBF clásico}
\subsection{Qué se hiperreduce y qué no}
Con ECM clásico:
\begin{itemize}[leftmargin=1.5em]
\item \textbf{Residual/Tangente}: se ensamblan en soporte ECM fijo.
\item \textbf{Homogenización}: se hace con pesos fijos ($w_\varepsilon,w_\sigma$)
o con full-mesh, según configuración.
\item \textbf{Incógnita de Newton}: sigue siendo $q_m$.
\end{itemize}

\subsection{Fórmula operativa}
\[
r_r^{hprom}(q_m)\approx
\left(\frac{\partial u_f}{\partial q_m}\right)^\top
\sum_{e\in\mathcal Z_{res}} w^{res}_e\,r_e(u(q_m)).
\]
\[
K_r^{hprom}(q_m)\approx
\left(\frac{\partial u_f}{\partial q_m}\right)^\top
\left(\sum_{e\in\mathcal Z_{res}} w^{res}_e\,K_e(u(q_m))\right)
\left(\frac{\partial u_f}{\partial q_m}\right).
\]
Homogenización (ejemplo en stress):
\[
\bar{\sigma}(q_m)\approx\sum_{e\in\mathcal Z_{\sigma}} w^{\sigma}_e\,\bar{\sigma}_e(q_m).
\]

\subsection{Recordatorio de $C_m$ y $C_s$ (porque aparece en el código)}
Aunque en online resuelves en $q_m$, el código puede reconstruir:
\[
q_{pod}=C_m(A_m q_m)+C_s q_s.
\]
\begin{itemize}[leftmargin=1.5em]
\item $C_m$: mezcla la parte master hacia coordenada POD.
\item $C_s$: mezcla la parte slave hacia coordenada POD.
\item Ambos salen de un ajuste LS offline (Stage2b), no se ``inventan'' online.
\end{itemize}

\section{Stage8a: dataset estructurado MAW}
\subsection{Qué aporta Stage8a frente a Stage6a}
Stage6a construye un problema ECM \textbf{global}. Stage8a lo reordena por
\textbf{nodo estructurado del manifold} para poder imponer restricciones locales
$A_jw_j=b_j$ (una por cada estado estructurado $j$).

\subsection{Qué se guarda exactamente}
Con $N_n=451$ nodos estructurados y $n_e=990$ elementos:
\[
Q_{res}\in\mathbb{R}^{(n_qN_n)\times n_e},
\qquad
b_{res}\in\mathbb{R}^{n_qN_n}.
\]
En tu caso típico ($n_q=13$):
\[
Q_{res}\in\mathbb{R}^{(13\cdot451)\times990}.
\]

Si activas \texttt{--include-homogenization 1}, además:
\[
C_{hom}\in\mathbb{R}^{(6N_n)\times n_e},
\qquad
b_{hom}\in\mathbb{R}^{6N_n}.
\]

\subsection{Bloques por nodo (la pieza que luego usa Stage8b)}
Para cada nodo estructurado $j\in\{1,\dots,N_n\}$:
\begin{itemize}[leftmargin=1.5em]
\item bloque residual local:
\[
Q_{res}^{(j)}\in\mathbb{R}^{n_q\times n_e};
\]
\item bloque de homogenización local (si está habilitado):
\[
C_{hom}^{(j)}\in\mathbb{R}^{6\times n_e}.
\]
\item coordenada de estado asociada:
\[
q_m^{(j)}\in\mathbb{R}^{d_m}.
\]
\end{itemize}

\section{Stage8b: MAW-ECM (Fase 1 actual, \textbf{sin grafo})}
\subsection{Qué problema resolvemos hoy}
En la fase 1 actual no estamos usando acoplamiento por grafo en el pruning
(\texttt{--use-global-graph-2ndstage 0}).
El objetivo es reducir soporte manteniendo factibilidad local exacta:
\[
A_jw^{(j)}=b_j,\qquad w^{(j)}\ge 0,\qquad j=1,\dots,N_n.
\]

\subsection{Variables y tamaños}
\begin{itemize}[leftmargin=1.5em]
\item Soporte candidato inicial:
\[
\mathcal Z_{ini},\qquad n_c:=|\mathcal Z_{ini}|.
\]
\item Pesos adaptativos por nodo estructurado:
\[
W_{adapt}=
\begin{bmatrix}
\vert & \vert &  & \vert\\
w^{(1)} & w^{(2)} & \cdots & w^{(N_n)}\\
\vert & \vert &  & \vert
\end{bmatrix}
\in\mathbb{R}^{n_c\times N_n}.
\]
\item Conjunto activo de filas/candidatos:
\[
I_{activo}\subseteq\{1,\dots,n_c\}.
\]
\textit{Nota de lectura de código:} en scripts MATLAB/Python este conjunto suele aparecer
como \texttt{iCAND}.
\end{itemize}

\subsection{Paso 0: bootstrap ECM fijo}
Primero se genera una regla fija inicial factible:
\[
(\mathcal Z_{ini},w_{ini}).
\]
Según el modo:
\begin{itemize}[leftmargin=1.5em]
\item \texttt{maw-mode=res\_only}: bootstrap desde $Q_{res}$.
\item \texttt{maw-mode=single\_res\_hom}: bootstrap sobre bloque combinado
$[Q_{res};C_\varepsilon;C_\sigma]$.
\end{itemize}
Luego se inicializa:
\[
w^{(j)}\equiv w_{ini}\quad \forall j.
\]

\subsection{Paso 1: construir restricciones locales}
\subsubsection*{Modo \texttt{res\_only}}
\[
A_j = Q_{res}^{(j)}(:,\mathcal Z_{ini})\in\mathbb{R}^{n_q\times n_c},
\qquad
b_j = A_jw_{ini}\in\mathbb{R}^{n_q}.
\]

\subsubsection*{Modo \texttt{single\_res\_hom}}
\[
A_j=
\begin{bmatrix}
Q_{res}^{(j)}(:,\mathcal Z_{ini})\\
C_{hom}^{(j)}(:,\mathcal Z_{ini})
\end{bmatrix}
\in\mathbb{R}^{(n_q+6)\times n_c},
\qquad
b_j=A_jw_{ini}.
\]

\subsubsection*{Restricción opcional de suma de pesos}
Si activas \texttt{--enforce-sum-weights 1}, a cada bloque local se le agrega una fila:
\[
\mathbf{1}^\top w^{(j)} = s_{target}.
\]
Es decir:
\[
\widetilde A_j=
\begin{bmatrix}
A_j\\
\mathbf{1}^\top
\end{bmatrix},
\qquad
\widetilde b_j=
\begin{bmatrix}
b_j\\
s_{target}
\end{bmatrix}.
\]
Esto aumenta rigidez del problema local y puede imponer un piso en el soporte mínimo alcanzable.

\subsection{Paso 2: ciclo de pruning (estilo Joaquín, primera fase)}
\subsubsection*{Paso 2 explicado con tus números (153 candidatos, 451 nodos)}
En tu caso:
\[
n_c=153,\quad N_n=451,\quad n_q=13.
\]
Al terminar el Paso 1 tienes, para cada nodo $j$:
\[
A_j\in\mathbb{R}^{13\times 153},\qquad b_j\in\mathbb{R}^{13},
\]
y una inicialización global:
\[
W_{adapt}\in\mathbb{R}^{153\times451},
\quad
W_{adapt}(:,j)=w_{ini}\;\;\forall j.
\]

El objetivo del Paso 2 es quitar columnas/candidatos de forma \textbf{común} para todos
los nodos, sin perder factibilidad local.

Iteración típica del pruning:
\begin{enumerate}[leftmargin=1.5em]
\item Estado actual: $k=|I_{activo}|$ candidatos vivos (al inicio $k=153$).
\item Se propone eliminar un candidato $p\in I_{activo}$.
\item Queda conjunto tentativo:
\[
I_{keep}=I_{activo}\setminus\{p\},
\quad |I_{keep}|=k-1.
\]
\item Para cada nodo $j=1,\dots,451$ se trabaja con:
\[
A_{j,keep}=A_j(:,I_{keep})\in\mathbb{R}^{13\times(k-1)}.
\]
\item Para ese nodo, con peso previo
\[
w_{j,before}=W_{adapt}(I_{keep},j)\in\mathbb{R}^{k-1},
\]
se fuerza de nuevo la ecuación local:
\[
A_{j,keep}\,w_{j,after}=b_j.
\]
\item Se busca el cambio mínimo:
\[
w_{j,after}=w_{j,before}+\Delta w_j,\quad
A_{j,keep}\Delta w_j=b_j-A_{j,keep}w_{j,before}.
\]
En código (NOENF) se resuelve con:
\[
\Delta w_j=\texttt{np.linalg.lstsq}(A_{j,keep},rhs).
\]
\item Si para \textbf{algún} nodo falla rango o sale peso negativo, ese candidato $p$
se rechaza.
\item Si para \textbf{todos} los nodos sale factible, se acepta:
\[
I_{activo}\leftarrow I_{keep},\qquad
W_{adapt}(p,:)=0.
\]
\end{enumerate}

Esto se repite hasta no poder eliminar más o hasta alcanzar $n_{stop}$.

\subsubsection*{Qué significa ``candidato común''}
Es clave: no estás quitando un elemento para un nodo sí y otro no.
Cuando eliminas $p$, lo eliminas para \textbf{todas} las 451 columnas de $W_{adapt}$.
Por eso el soporte final es común y luego se puede usar en un único HROM mdpa.

\subsubsection*{Por qué puede estancarse}
Aunque $k-1\ge13$, puede pasar que:
\begin{itemize}[leftmargin=1.5em]
\item para algunos nodos $A_{j,keep}$ pierda rango efectivo;
\item o la solución exacta local requiera entradas negativas.
\end{itemize}
Cuando eso pasa para muchos candidatos, el algoritmo ya no puede seguir podando.

\subsubsection*{NOENF vs NOENFe en una frase}
\begin{itemize}[leftmargin=1.5em]
\item NOENF: intenta actualización rápida y rechaza si salen negativos.
\item NOENFe: activa enforcement explícito de no-negatividad por active-set local.
\end{itemize}

Si NOENF no puede reducir más, NOENFe intenta rescatar más eliminaciones factibles.

\subsubsection*{Ejemplo numérico completo (matrices explícitas)}
Para ver el mecanismo de verdad, toma un caso pequeño:
\[
n_q=2,\quad n_c=4,\quad N_n=3.
\]
Soporte inicial:
\[
\mathcal Z_{ini}=\{1,2,3,4\}.
\]
Peso fijo inicial:
\[
w_{ini}=
\begin{bmatrix}
0.5\\0.7\\0.6\\0.4
\end{bmatrix}.
\]

Definimos tres nodos estructurados con:
\[
A_1=
\begin{bmatrix}
1&0&1&1\\
0&1&1&1
\end{bmatrix},\quad
A_2=
\begin{bmatrix}
1&0&2&2\\
0&1&1&1
\end{bmatrix},\quad
A_3=
\begin{bmatrix}
2&0&1&1\\
0&2&1&1
\end{bmatrix}.
\]
Y:
\[
b_j=A_jw_{ini}.
\]
Entonces:
\[
b_1=
\begin{bmatrix}
1.5\\1.7
\end{bmatrix},\quad
b_2=
\begin{bmatrix}
2.5\\1.7
\end{bmatrix},\quad
b_3=
\begin{bmatrix}
2.0\\2.4
\end{bmatrix}.
\]

\paragraph{Paso 0 (ejemplo pequeño): cómo se elige el candidato}
Aquí la decisión \textbf{no} se toma con un único $w_{ini}$, sino con la matriz
adaptativa actual:
\[
W_{adapt}^{(t)}\in\mathbb{R}^{n_c\times N_n}.
\]
En este ejemplo:
\[
n_c=4,\qquad N_n=3.
\]
Al inicio:
\[
W_{adapt}^{(0)}=
\begin{bmatrix}
0.5&0.5&0.5\\
0.7&0.7&0.7\\
0.6&0.6&0.6\\
0.4&0.4&0.4
\end{bmatrix},
\]
porque en $t=0$ todas las columnas son $w_{ini}$.

Se calcula la actividad por candidato:
\[
a_i^{(t)}=\sum_{j=1}^{N_n}W_{adapt}^{(t)}(i,j).
\]
Entonces:
\[
a^{(0)}=
\begin{bmatrix}
1.5\\2.1\\1.8\\1.2
\end{bmatrix},
\]
y el orden ascendente es:
\[
4\;<\;1\;<\;3\;<\;2.
\]
Por eso el \textbf{primer intento real} es eliminar el candidato 4.

\paragraph{Intento A: eliminar candidato 4 (se acepta)}
Queda:
\[
I_{keep}=\{1,2,3\}.
\]
Matrices literales de este intento:
\[
A_{1,keep}=A_1(:,\{1,2,3\})=
\begin{bmatrix}
1&0&1\\
0&1&1
\end{bmatrix},
\]
\[
A_{2,keep}=A_2(:,\{1,2,3\})=
\begin{bmatrix}
1&0&2\\
0&1&1
\end{bmatrix},
\]
\[
A_{3,keep}=A_3(:,\{1,2,3\})=
\begin{bmatrix}
2&0&1\\
0&2&1
\end{bmatrix}.
\]
Pesos ``before'' literales:
\[
w_{1,before}=w_{2,before}=w_{3,before}=
\begin{bmatrix}
0.5\\0.7\\0.6
\end{bmatrix}.
\]
Definiciones (explícitas) para este intento:
\begin{itemize}[leftmargin=1.5em]
\item $A_{j,keep}\in\mathbb{R}^{n_q\times |I_{keep}|}$:
submatriz de $A_j$ tomando solo las columnas de índices en $I_{keep}$.
En este ejemplo, $n_q=2$ y $|I_{keep}|=3$, por tanto
$A_{j,keep}\in\mathbb{R}^{2\times 3}$.
\item $w_{j,before}\in\mathbb{R}^{|I_{keep}|}$:
vector de pesos del nodo $j$ antes de intentar eliminar el candidato,
restringido a las filas activas $I_{keep}$.
\item $\Delta w_j\in\mathbb{R}^{|I_{keep}|}$:
incremento que corrige $w_{j,before}$ para mantener exactitud local.
\item $w_{j,after}=w_{j,before}+\Delta w_j$:
pesos locales tras el intento.
\end{itemize}

En cada nodo:
\[
A_{j,keep}\Delta w_j=b_j-A_{j,keep}w_{j,before},
\qquad
w_{j,after}=w_{j,before}+\Delta w_j.
\]
Definimos para hacerlo ultra explícito:
\[
y_{j,before}:=A_{j,keep}w_{j,before},
\qquad
y_{j,after}:=A_{j,keep}w_{j,after}.
\]
Con $w_{j,before}=[0.5,0.7,0.6]^\top$, la solución de mínimo cambio da:
\[
\Delta w_1=
\begin{bmatrix}
0.1333\\0.1333\\0.2667
\end{bmatrix},\quad
\Delta w_2=
\begin{bmatrix}
0.1333\\0.0667\\0.3333
\end{bmatrix},\quad
\Delta w_3=
\begin{bmatrix}
0.1333\\0.1333\\0.1333
\end{bmatrix}.
\]
Por tanto:
\[
w_{1,after}=
\begin{bmatrix}
0.6333\\0.8333\\0.8667
\end{bmatrix},\quad
w_{2,after}=
\begin{bmatrix}
0.6333\\0.7667\\0.9333
\end{bmatrix},\quad
w_{3,after}=
\begin{bmatrix}
0.6333\\0.8333\\0.7333
\end{bmatrix}.
\]
Todos son no-negativos y satisfacen $A_{j,keep}w_{j,after}=b_j$, así que se acepta
la eliminación del 4.
\[
y_{1,before}=
\begin{bmatrix}
1.1\\1.3
\end{bmatrix},\;
y_{1,after}=
\begin{bmatrix}
1.5\\1.7
\end{bmatrix},\;
y_{2,before}=
\begin{bmatrix}
1.7\\1.3
\end{bmatrix},\;
y_{2,after}=
\begin{bmatrix}
2.5\\1.7
\end{bmatrix},\;
y_{3,before}=
\begin{bmatrix}
2.0\\2.7
\end{bmatrix},\;
y_{3,after}=
\begin{bmatrix}
2.0\\2.4
\end{bmatrix}.
\]

La nueva matriz adaptativa queda:
\[
W_{adapt}^{(1)}=
\begin{bmatrix}
0.6333&0.6333&0.6333\\
0.8333&0.7667&0.8333\\
0.8667&0.9333&0.7333\\
0&0&0
\end{bmatrix}.
\]
Ahora:
\[
a^{(1)}=
\begin{bmatrix}
1.900\\2.4333\\2.5333\\0
\end{bmatrix},
\]
de modo que, sobre $I_{positivo}^{(1)}=\{1,2,3\}$, el siguiente candidato real es
el 1 (no el 2).

\paragraph{Mismo Paso 1 pero imponiendo \texorpdfstring{$\mathbf 1^\top w=S_{target}$}{1^T w = S_target}}
Ahora repetimos \textbf{solo} el intento A (eliminar 4), pero añadiendo:
\[
\mathbf 1^\top w_{j,after}=S_{target}.
\]
Para este ejemplo tomamos:
\[
S_{target}=2.2=\mathbf 1^\top
\begin{bmatrix}
0.5\\0.7\\0.6\\0.4
\end{bmatrix},
\]
que es la suma total inicial de pesos en el soporte de 4 candidatos.

Entonces, para cada nodo:
\[
\widetilde A_{j,keep}w_{j,after}=\widetilde b_j,\qquad
\widetilde A_{j,keep}=
\begin{bmatrix}
A_{j,keep}\\
1\;1\;1
\end{bmatrix}.
\]

Nodo 1:
\[
\begin{bmatrix}
1&0&1\\
0&1&1\\
1&1&1
\end{bmatrix}
\begin{bmatrix}
w_1\\w_2\\w_3
\end{bmatrix}
=
\begin{bmatrix}
1.5\\1.7\\2.2
\end{bmatrix}
\Rightarrow
w_{1,after}^{sum}=
\begin{bmatrix}
0.5\\0.7\\1.0
\end{bmatrix}.
\]
Nodo 2:
\[
\begin{bmatrix}
1&0&2\\
0&1&1\\
1&1&1
\end{bmatrix}
\begin{bmatrix}
w_1\\w_2\\w_3
\end{bmatrix}
=
\begin{bmatrix}
2.5\\1.7\\2.2
\end{bmatrix}
\Rightarrow
w_{2,after}^{sum}=
\begin{bmatrix}
0.5\\0.7\\1.0
\end{bmatrix}.
\]
Nodo 3:
\[
\begin{bmatrix}
2&0&1\\
0&2&1\\
1&1&1
\end{bmatrix}
\begin{bmatrix}
w_1\\w_2\\w_3
\end{bmatrix}
=
\begin{bmatrix}
2.0\\2.4\\2.2
\end{bmatrix}.
\]
La tercera ecuación aquí es combinación lineal de las dos primeras, y el
mínimo cambio da:
\[
w_{3,after}^{sum}=
\begin{bmatrix}
0.6333\\0.8333\\0.7333
\end{bmatrix}.
\]

Con esto:
\[
W_{adapt}^{(1,sum)}=
\begin{bmatrix}
0.5000&0.5000&0.6333\\
0.7000&0.7000&0.8333\\
1.0000&1.0000&0.7333\\
0&0&0
\end{bmatrix}.
\]
Observa que para cada nodo se cumple:
\[
\mathbf 1^\top w_{j,after}^{sum}=2.2.
\]

\textbf{Comparación directa del Paso 1:}
\begin{itemize}[leftmargin=1.5em]
\item Sin restricción de suma:
\[
W_{adapt}^{(1)}=
\begin{bmatrix}
0.6333&0.6333&0.6333\\
0.8333&0.7667&0.8333\\
0.8667&0.9333&0.7333\\
0&0&0
\end{bmatrix}.
\]
\item Con $\mathbf 1^\top w=S_{target}$:
\[
W_{adapt}^{(1,sum)}=
\begin{bmatrix}
0.5000&0.5000&0.6333\\
0.7000&0.7000&0.8333\\
1.0000&1.0000&0.7333\\
0&0&0
\end{bmatrix}.
\]
\end{itemize}

Desde aquí, para no mezclar ramas, continuamos el ejemplo principal por la rama
\textbf{sin} restricción de suma (la que ya veníamos usando).

\paragraph{Intento B: eliminar candidato 1 (se acepta)}
Ahora:
\[
I_{keep}=\{2,3\}.
\]
Definiciones (mismas ideas, nuevo $I_{keep}$):
\begin{itemize}[leftmargin=1.5em]
\item $A_{j,keep}$: submatriz de $A_j$ con columnas $\{2,3\}$.
Aquí $A_{j,keep}\in\mathbb{R}^{2\times 2}$.
\item $w_{j,before}\in\mathbb{R}^{2}$:
columna $j$ de $W_{adapt}^{(1)}$ restringida a filas $\{2,3\}$.
\item $\Delta w_j\in\mathbb{R}^{2}$ y
$w_{j,after}=w_{j,before}+\Delta w_j$.
\end{itemize}

Para este intento, en cada nodo trabajamos con:
\[
A_{j,keep}\Delta w_j
=
b_j-A_{j,keep}w_{j,before},
\qquad
w_{j,after}=w_{j,before}+\Delta w_j.
\]
Y otra vez:
\[
y_{j,before}:=A_{j,keep}w_{j,before},
\qquad
y_{j,after}:=A_{j,keep}w_{j,after}.
\]
Aquí $w_{j,before}$ se toma de $W_{adapt}^{(1)}$ en las filas $\{2,3\}$.

Nodo 1 ($j=1$):
\[
A_{1,keep}=
\begin{bmatrix}
0&1\\
1&1
\end{bmatrix}
\begin{bmatrix}
w_2\\w_3
\end{bmatrix}
=
\begin{bmatrix}
1.5\\1.7
\end{bmatrix}
\]
\[
w_{1,before}=
\begin{bmatrix}
0.8333\\0.8667
\end{bmatrix},
\quad
y_{1,before}=A_{1,keep}w_{1,before}=
\begin{bmatrix}
0.8667\\1.7000
\end{bmatrix},
\]
\[
RHS_1=b_1-A_{1,keep}w_{1,before}
=
\begin{bmatrix}
1.5\\1.7
\end{bmatrix}
-
\begin{bmatrix}
0.8667\\1.7000
\end{bmatrix}
=
\begin{bmatrix}
0.6333\\0
\end{bmatrix}.
\]
Resolviendo:
\[
\Delta w_1=
\begin{bmatrix}
-0.6333\\0.6333
\end{bmatrix},
\qquad
w_{1,after}=
\begin{bmatrix}
0.2000\\1.5000
\end{bmatrix}.
\]
Chequeo:
\[
y_{1,after}=A_{1,keep}w_{1,after}
=
\begin{bmatrix}
1.5\\1.7
\end{bmatrix}
=b_1.
\]

Nodo 2 ($j=2$):
\[
A_{2,keep}=
\begin{bmatrix}
0&2\\
1&1
\end{bmatrix}
\begin{bmatrix}
w_2\\w_3
\end{bmatrix}
=
\begin{bmatrix}
2.5\\1.7
\end{bmatrix}
\]
\[
w_{2,before}=
\begin{bmatrix}
0.7667\\0.9333
\end{bmatrix},
\quad
y_{2,before}=A_{2,keep}w_{2,before}=
\begin{bmatrix}
1.8666\\1.7000
\end{bmatrix},
\]
\[
RHS_2=b_2-A_{2,keep}w_{2,before}
=
\begin{bmatrix}
2.5\\1.7
\end{bmatrix}
-
\begin{bmatrix}
1.8666\\1.7000
\end{bmatrix}
=
\begin{bmatrix}
0.6334\\0
\end{bmatrix}.
\]
Resolviendo:
\[
\Delta w_2=
\begin{bmatrix}
-0.3167\\0.3167
\end{bmatrix},
\qquad
w_{2,after}=
\begin{bmatrix}
0.4500\\1.2500
\end{bmatrix}.
\]
Chequeo:
\[
y_{2,after}=A_{2,keep}w_{2,after}
=
\begin{bmatrix}
2.5\\1.7
\end{bmatrix}
=b_2.
\]

Nodo 3 ($j=3$):
\[
A_{3,keep}=
\begin{bmatrix}
0&1\\
2&1
\end{bmatrix}
\begin{bmatrix}
w_2\\w_3
\end{bmatrix}
=
\begin{bmatrix}
2.0\\2.4
\end{bmatrix}
\]
\[
w_{3,before}=
\begin{bmatrix}
0.8333\\0.7333
\end{bmatrix},
\quad
y_{3,before}=A_{3,keep}w_{3,before}=
\begin{bmatrix}
0.7333\\2.3999
\end{bmatrix},
\]
\[
RHS_3=b_3-A_{3,keep}w_{3,before}
=
\begin{bmatrix}
2.0\\2.4
\end{bmatrix}
-
\begin{bmatrix}
0.7333\\2.3999
\end{bmatrix}
\approx
\begin{bmatrix}
1.2667\\0
\end{bmatrix}.
\]
Resolviendo:
\[
\Delta w_3=
\begin{bmatrix}
-0.6333\\1.2667
\end{bmatrix},
\qquad
w_{3,after}=
\begin{bmatrix}
0.2000\\2.0000
\end{bmatrix}.
\]
Chequeo:
\[
y_{3,after}=A_{3,keep}w_{3,after}
=
\begin{bmatrix}
2.0\\2.4
\end{bmatrix}
=b_3.
\]

Todos los $w_{j,after}$ son no-negativos, así que se acepta eliminar el candidato 1.
La nueva matriz adaptativa es:
\[
W_{adapt}^{(2)}=
\begin{bmatrix}
0&0&0\\
0.2000&0.4500&0.2000\\
1.5000&1.2500&2.0000\\
0&0&0
\end{bmatrix}.
\]
Luego:
\[
a^{(2)}=
\begin{bmatrix}
0\\0.8500\\4.7500\\0
\end{bmatrix},
\]
y en $I_{positivo}^{(2)}=\{2,3\}$ el siguiente candidato natural es el 2.

\paragraph{Intento C: eliminar candidato 2 (se rechaza)}
Tras el intento B:
\[
I_{activo}=\{2,3\},
\]
y el siguiente en ranking es el 2. Si lo quitas:
\[
I_{keep}=\{3\}.
\]
Definiciones en este caso:
\begin{itemize}[leftmargin=1.5em]
\item $A_{j,keep}$: submatriz de $A_j$ con la única columna 3
($A_{j,keep}\in\mathbb{R}^{2\times 1}$).
\item $w_{j,before}\in\mathbb{R}^{1}$: peso actual del candidato 3 en el nodo $j$.
\item $\Delta w_j\in\mathbb{R}^{1}$: único incremento escalar a resolver.
\end{itemize}

Usando $W_{adapt}^{(2)}$, para nodo 2:
\[
w_{2,before}=
\begin{bmatrix}
1.2500
\end{bmatrix},
\qquad
A_{2,keep}=
\begin{bmatrix}
2\\1
\end{bmatrix},
\qquad
b_2=
\begin{bmatrix}
2.5\\1.7
\end{bmatrix}.
\]
\[
y_{2,before}=A_{2,keep}w_{2,before}
=
\begin{bmatrix}
2.5\\1.25
\end{bmatrix},
\qquad
RHS_2=
b_2-A_{2,keep}w_{2,before}
=
\begin{bmatrix}
0\\0.45
\end{bmatrix}.
\]
Con incremento escalar $\Delta w_2$:
\[
2(1.25+\Delta w_2)=2.5,\qquad
1.25+\Delta w_2=1.7.
\]
Equivalentemente:
\[
\Delta w_2=0,\qquad
\Delta w_2=0.45.
\]
Incompatible. No existe un único $\Delta w_2$ que satisfaga ambas ecuaciones exactas.
Por eso este candidato se rechaza.

La matriz queda sin cambios:
\[
W_{adapt}^{(3)}=W_{adapt}^{(2)}.
\]

\paragraph{Variante: imponer conservación de suma de pesos}
Si quieres imponer, en cada nodo $j$, que la suma de pesos se conserve, no cambias
la lógica base: solo añades una ecuación lineal extra al sistema local.

Sistema base local:
\[
A_{j,keep}w_{j,after}=b_j.
\]
Sistema aumentado:
\[
\widetilde A_{j,keep}w_{j,after}=\widetilde b_j,
\qquad
\widetilde A_{j,keep}=
\begin{bmatrix}
A_{j,keep}\\
c^\top
\end{bmatrix},
\quad
\widetilde b_j=
\begin{bmatrix}
b_j\\
c^\top w_{j,ref}
\end{bmatrix}.
\]

\textbf{Caso correcto para conservar suma:}
\[
c=\mathbf 1.
\]
Entonces:
\[
\mathbf 1^\top w_{j,after}=\mathbf 1^\top w_{j,ref}.
\]
Si en tu caso real quieres preservar $990$, usas:
\[
c=\mathbf 1_{153},\qquad
\mathbf 1_{153}^\top w_{j,ref}=990.
\]

\textbf{No es lo mismo usar $c=w_{ini}$:}
\[
w_{ini}^\top w_{j,after}=w_{ini}^\top w_{j,ref},
\]
eso conserva un producto ponderado, no la suma.

\paragraph{Mini-cálculo explícito (ejemplo pequeño)}
Tomamos el nodo 2 del intento A con:
\[
A_{2,keep}=
\begin{bmatrix}
1&0&2\\
0&1&1
\end{bmatrix},\quad
b_2=
\begin{bmatrix}
2.5\\1.7
\end{bmatrix},\quad
w_{2,ref}=w_{2,before}=
\begin{bmatrix}
0.5\\0.7\\0.6
\end{bmatrix}.
\]

Si impones conservación de suma con $c=\mathbf 1_3$:
\[
\widetilde A_{2,keep}=
\begin{bmatrix}
1&0&2\\
0&1&1\\
1&1&1
\end{bmatrix},\qquad
\widetilde b_2=
\begin{bmatrix}
2.5\\1.7\\1.8
\end{bmatrix}
\]
(porque $1.8=0.5+0.7+0.6$).

Resolviendo:
\[
w_{2,after}=
\begin{bmatrix}
0.1\\0.5\\1.2
\end{bmatrix},
\qquad
\Delta w_2=
\begin{bmatrix}
-0.4\\-0.2\\0.6
\end{bmatrix}.
\]
Chequeos:
\[
A_{2,keep}w_{2,after}=
\begin{bmatrix}
2.5\\1.7
\end{bmatrix}=b_2,
\qquad
\mathbf 1^\top w_{2,after}=1.8=\mathbf 1^\top w_{2,ref}.
\]

Si en cambio usas $c=w_{2,ref}=[0.5,0.7,0.6]^\top$:
\[
w_{2,after}\approx
\begin{bmatrix}
0.0636\\0.4818\\1.2182
\end{bmatrix},
\]
y:
\[
\mathbf 1^\top w_{2,after}\approx 1.7636 \neq 1.8.
\]
Con esto ves por qué, si el objetivo es conservar \emph{suma}, la fila extra debe
ser $\mathbf 1^\top$ (o la medida geométrica que realmente quieras conservar).

\textbf{Regla global de decisión (en el ejemplo pequeño y en el caso general):}
\begin{itemize}[leftmargin=1.5em]
\item El ranking se hace con $W_{adapt}^{(t)}$ actual, no con $w_{ini}$ fijo.
\item Se acepta un candidato solo si \emph{todos} los nodos son factibles.
\item Si falla un nodo (incompatibilidad, rango o negatividad), el candidato se rechaza.
\end{itemize}

\paragraph{Qué aprendiste aquí}
\begin{itemize}[leftmargin=1.5em]
\item $\Delta w_j$ es vector para cada nodo $j$.
\item Si apilas todos los nodos, obtienes un bloque de cambios
\[
\Delta W=
\begin{bmatrix}
\vert&\vert&\vert\\
\Delta w_1&\Delta w_2&\Delta w_3\\
\vert&\vert&\vert
\end{bmatrix}.
\]
\item La decisión aceptar/rechazar un candidato es global:
si falla un solo nodo, se rechaza ese candidato.
\end{itemize}

Sea
\[
a_i:=\sum_{j=1}^{N_n}W_{adapt}(i,j),
\]
que interpretamos como \textbf{actividad total del candidato $i$} en todos los nodos
del manifold.
Se define:
\[
I_{positivo}=\{i\in I_{activo}:a_i>0\},
\]
y se ordenan candidatos por $a_i$ ascendente (se intenta eliminar primero los ``menos usados'').
\textit{Nota de lectura de código:} $a_i$ suele aparecer como \texttt{RRR\_i} y
$I_{positivo}$ como \texttt{ilocPOS}.

\subsubsection*{Etapa 2.1: NOENF (sin enforcement explícito de positividad)}
Para cada candidato propuesto:
\begin{enumerate}[leftmargin=1.5em]
\item quitar temporalmente la fila candidata;
\item para cada nodo $j$, resolver actualización de mínimo cambio:
\[
A_{j,keep}\,\Delta w_j=b_j-A_{j,keep}w_{j,before},
\]
\[
w_{j,after}=w_{j,before}+\Delta w_j;
\]
\item si hay rank deficiency o algún peso negativo ($<-\texttt{tol\_neg}$), se rechaza ese candidato;
\item si todos los nodos son factibles, se acepta la eliminación.
\end{enumerate}
Esto corresponde a una \textbf{pasada de eliminación sin enforcement explícito de positividad}
(\texttt{Loop\_MAWecmNOENF\_cl}-like).

\subsubsection*{Etapa 2.2: NOENFe local (positividad explícita, sin grafo)}
Si NOENF no puede eliminar más, se pasa a active-set local
(\texttt{Loop\_MAWecmNOENFe}-like), todavía \textbf{sin} acoplar nodos con grafo:
\begin{itemize}[leftmargin=1.5em]
\item se impone no negatividad explícita durante la actualización local;
\item se acepta el primer candidato factible según el orden.
\end{itemize}

\subsubsection*{Qué sistema lineal se resuelve en NOENF (exactamente)}
Para un candidato a eliminar y un nodo $j$:
\[
A_{j,keep}\in\mathbb{R}^{m_j\times r_{keep}},
\qquad
w_{j,before}\in\mathbb{R}^{r_{keep}}.
\]
Se resuelve:
\[
A_{j,keep}\,\Delta w_j = b_j - A_{j,keep}w_{j,before}.
\]
En código:
\[
\Delta w_j=\texttt{np.linalg.lstsq}(A_{j,keep},\;rhs).
\]
Luego:
\[
w_{j,after}=w_{j,before}+\Delta w_j.
\]
Se rechaza el candidato si:
\begin{itemize}[leftmargin=1.5em]
\item \texttt{matrix\_rank}($A_{j,keep}$) es insuficiente;
\item existe entrada negativa por debajo de \texttt{-tol\_neg}.
\end{itemize}

\subsubsection*{Qué sistema se resuelve en NOENFe local (active-set)}
En la implementación actual (\texttt{mawecm\_pruning.py}),
por cada nodo $j$ con conjunto libre $F_j$:
\[
A_{j,F_j}\in\mathbb{R}^{m_j\times |F_j|}.
\]
Se calcula una solución particular:
\[
w_{p,j}=A_{j,F_j}^\top\left(A_{j,F_j}A_{j,F_j}^\top\right)^{-1}b_j.
\]
En código no se forma la inversa explícita, se usa:
\[
\texttt{np.linalg.solve}(G_m,b_j),\quad G_m=A_{j,F_j}A_{j,F_j}^\top,
\]
y si falla:
\[
\texttt{np.linalg.lstsq}(G_m,b_j).
\]

Después se construye base de null-space $N_j$ (vía SVD) y se resuelve el sistema reducido:
\[
\left(N^\top HN\right)y=N^\top\left(g-Hz_p\right),
\qquad
z=z_p+Ny.
\]
En código:
\begin{itemize}[leftmargin=1.5em]
\item intento 1: \texttt{scipy.sparse.linalg.spsolve};
\item fallback: \texttt{np.linalg.lstsq}.
\end{itemize}

En la fase 1 sin grafo:
\[
\alpha_{smooth}=0,\quad K_{graph}=0\quad\Rightarrow\quad H=I,
\]
así que este paso se simplifica bastante.

\subsubsection*{Condición de parada}
Se detiene cuando:
\begin{itemize}[leftmargin=1.5em]
\item $|I_{activo}|\le n_{stop}$ (definido por \texttt{--maw-min-support-size} o auto), o
\item no hay más eliminaciones factibles.
\end{itemize}

\subsection{Paso 3: salida del pruning}
El resultado principal es:
\[
\mathcal Z_{red}\subseteq\mathcal Z_{ini},
\qquad
W_{train}\in\mathbb{R}^{|\mathcal Z_{red}|\times N_n}.
\]
Además se valida:
\[
\max_j\frac{\|A_jw^{(j)}-b_j\|_2}{\|b_j\|_2}\le \texttt{strict\_constraint\_rel\_tol},
\qquad
\min_{i,j}W_{train}(i,j)\ge -\texttt{strict\_negative\_tol}.
\]

\subsection{Paso 4: RBF final de pesos adaptativos}
Con pares de entrenamiento:
\[
\big\{q_m^{(j)},\,w^{(j)}\big\}_{j=1}^{N_n},
\]
se ajusta RBF multi-output:
\[
q_m\mapsto w(q_m)\in\mathbb{R}^{|\mathcal Z_{red}|}.
\]
En online:
\begin{itemize}[leftmargin=1.5em]
\item clipping no negativo (si está activo);
\item renormalización opcional a objetivo de suma (por defecto $\sum w_{ini}$).
\end{itemize}

\subsubsection*{Cómo se ajusta ese RBF (álgebra + función usada)}
En \texttt{mawecm\_rbf\_weights.py}:
\begin{itemize}[leftmargin=1.5em]
\item se arma $\Phi$ (kernel gaussiano entre muestras y centros);
\item si no hay parte polinómica:
\[
X=\texttt{np.linalg.lstsq}\!\left(\Phi^\top\Phi+\lambda I,\;\Phi^\top Y\right);
\]
\item si hay parte polinómica (constante o afín), se resuelve el sistema bloque:
\[
\begin{bmatrix}
\Phi^\top\Phi+\lambda I & \Phi^\top P\\
P^\top\Phi & P^\top P
\end{bmatrix}
\begin{bmatrix}
X_{\phi}\\ X_{poly}
\end{bmatrix}
=
\begin{bmatrix}
\Phi^\top Y\\ P^\top Y
\end{bmatrix},
\]
también con \texttt{np.linalg.lstsq}.
\end{itemize}

\subsubsection*{Homogenización por soporte MAW (modo \texttt{maw\_support\_nnls})}
Cuando se ajustan pesos de homogenización sobre soporte MAW reducido, el problema es:
\[
\min_{w\ge 0}\|Aw-b\|_2^2+\lambda\|w\|_2^2.
\]
En código \textbf{no} se usa \texttt{scipy.optimize.nnls} directo; se usa
\texttt{scipy.optimize.lsq\_linear} con cotas $(0,\infty)$ (NNLS acotado), con secuencia robusta:
\begin{enumerate}[leftmargin=1.5em]
\item \texttt{method='trf', lsq\_solver='lsmr'};
\item fallback \texttt{method='trf', lsq\_solver='exact'};
\item fallback \texttt{method='bvls'}.
\end{enumerate}
Si todas fallan, se lanza error explícito.

\subsection{Nota importante: ``sin grafo'' no significa ``sin Laplaciano en disco''}
El grafo estructurado se puede seguir guardando desde Stage0/Stage8a, pero en esta fase:
\[
\texttt{use-global-graph-2ndstage}=0
\]
y el pruning no acopla nodos por $K_{graph}$.
El acoplamiento por grafo queda para una segunda fase.

\section{Modos MAW actuales (implementación 2D)}
\subsection{\texttt{maw-mode=res\_only}}
\begin{itemize}[leftmargin=1.5em]
\item MAW adapta solo residual.
\item Homogenización en online puede venir de:
\texttt{full\_mesh}, \texttt{fixed\_ecm} o \texttt{maw\_support\_nnls}.
\item Es el modo más robusto para arrancar y depurar.
\end{itemize}

\subsection{\texttt{maw-mode=single\_res\_hom}}
\begin{itemize}[leftmargin=1.5em]
\item MAW se entrena con restricciones locales combinadas residual+hom.
\item En online, homogenización puede evaluarse con pesos MAW del mismo soporte.
\item Si además impones suma de pesos local, el soporte mínimo factible suele subir.
\end{itemize}

\subsection{Aclaración práctica sobre \texorpdfstring{$Z_{union}$}{Z_union}}
Cuando residual y homogenización usan soportes distintos:
\[
Z_{union}=Z_{res}\cup Z_\varepsilon\cup Z_\sigma.
\]
Por eso puedes pedir, por ejemplo, un tamaño objetivo en residual (p.ej. 100)
y terminar con un \textbf{union} mayor en la malla HROM final.
No es inconsistencia: son objetivos distintos compartiendo una malla reducida final.

\section{Stage8 online: qué se compara exactamente}
En \texttt{stage8\_test\_hprom\_mawecm\_rbf\_online.py} (o su variante GPR) se ejecuta la misma trayectoria para:
\begin{itemize}[leftmargin=1.5em]
\item FOM,
\item PROM,
\item HPROM-MAWECM.
\end{itemize}

Se reporta:
\begin{itemize}[leftmargin=1.5em]
\item error relativo en strain y stress macroscópicos;
\item error en coordenadas de estado:
\[
\|q_m^{HPROM}-q_m^{PROM}\|,\qquad
\|q_s^{HPROM}-q_s^{PROM}\|;
\]
\item tiempo total y \textit{speedup}.
\end{itemize}

Si activas malla HROM (\texttt{--hprom-use-hrom-mdpa 1}), el costo de ensamblaje cae;
si no, aunque haya pesos hiperreducidos, parte del costo de malla completa puede quedarse.

\section{Aclaraciones clave (las que más confunden)}
\subsection{``¿q\_pod se usa para resolver online?''}
En el flujo actual, la incógnita online es $q_m$.
$q_s$ viene del decoder RBF.
$q_{pod}$ puede guardarse como diagnóstico, pero no es la incógnita primaria del Newton.

\subsection{``¿Qué es $R_r$?''}
\[
R_r = \left(\frac{\partial u_f}{\partial q_m}\right)^\top R_f.
\]
Es el residual en el subespacio reducido.
El criterio principal de convergencia en PROM/HPROM está ligado a $\|R_r\|$.

\subsection{``¿Por qué aparece K\_old?''}
Se reutiliza la rigidez reducida del paso/iteración anterior como aceleración en la primera iteración correctora cuando aplica.

\subsection{``¿Qué significa single en ECM/MAW?''}
\begin{itemize}[leftmargin=1.5em]
\item En ECM clásico Stage6b: \texttt{single} = una sola selección/pesos para res+eps+sig.
\item En MAW Stage8b: \texttt{single\_res\_hom} = entrenamiento MAW usando restricciones de residual+hom juntas.
\end{itemize}

\section{Ejemplo mini de tamaños (con números tuyos)}
En una corrida típica:
\[
n_f=3924,\; r=13,\; d_m=2,\; d_s=11,\; n_e=990,\; N_n=451.
\]

\subsection{Decoder}
\[
\Phi_m\in\mathbb{R}^{3924\times 2},\;
A_m\in\mathbb{R}^{2\times 2},\;
\Phi_s\in\mathbb{R}^{3924\times 11}.
\]
Para un estado:
\[
q_m\in\mathbb{R}^{2}\;\Rightarrow\; q_s\in\mathbb{R}^{11}\;\Rightarrow\; u_f\in\mathbb{R}^{3924}.
\]

\subsection{MAW data}
\[
W_{train}\in\mathbb{R}^{|\mathcal{Z}_{red}|\times 451}.
\]
Si $|\mathcal{Z}_{red}|=172$, entonces:
\[
W_{train}\in\mathbb{R}^{172\times 451}.
\]

\section{Checklist práctico para depurar}
\begin{enumerate}[leftmargin=1.5em]
\item Verifica coherencia dimensional de \texttt{A\_m, C\_m, C\_s, phi\_m, phi\_s}.
\item Verifica que Stage4 esté en \texttt{rbf-input-space=q\_m}.
\item Verifica que \texttt{q\_m\_train} usado en Stage3 venga de ``consistent'' (no LS crudo).
\item En MAW-single, reconstruye Stage8a con \texttt{--include-homogenization 1}.
\item En MAW-single online usa \texttt{--hprom-homogenization-mode maw\_single}.
\item Si usas HROM mdpa, confirma mapeo full$\leftrightarrow$hrom en el \texttt{ecm\_weights\_all.npz}.
\end{enumerate}

\section{Conclusión}
Tu implementación actual ya tiene una estructura muy clara:
\begin{itemize}[leftmargin=1.5em]
\item una capa de \textbf{coordenadas reducidas coherentes} (Stage2a/2b);
\item un \textbf{decoder} $q_m\mapsto q_s$ por RBF (Stage4);
\item un PROM/HPROM con \textbf{Newton reducido en $q_m$};
\item ECM clásico para baseline;
\item MAW-ECM con pesos adaptativos en residual, y ahora también modo single res+hom.
\end{itemize}

\noindent
Si quieres, como siguiente documento puedo hacer uno aún más ``de pizarra'':
\begin{itemize}[leftmargin=1.5em]
\item solo 3 páginas;
\item un ejemplo numérico completo con $r=3$, $d_m=1$, $d_s=2$, $n_e=5$;
\item y resolver a mano un paso de Newton reducido + un paso de MAW pruning.
\end{itemize}

\end{document}
